% Source: physical PDF p. 405 (printed p. 382), beginning with the first
% paragraph below the Chapter 32 title and ending immediately before the
% Section 32.1 heading.
% Inventory: four linked citation markers ([1]--[4]); no numbered or
% unnumbered displays, footnotes, figures, tables, or typographic dividers.
% Uncertainties: none.

Ever since the ground-breaking work of Kaluza~\hyperref[ch32-ref-1]{[1]}
and Klein~\hyperref[ch32-ref-2]{[2]}, theorists have from time to time
tried to formulate a more nearly fundamental physical theory in spacetimes
of higher than four dimensions. This approach was revived in superstring
theories, which take their simplest form in 10 spacetime
dimensions.~\hyperref[ch32-ref-3]{[3]} More recently, it has been suggested
that the various versions of string theory may be unified in a theory known
as \(M\) theory, which in one limit is approximately described by
supergravity in 11 spacetime dimensions.~\hyperref[ch32-ref-4]{[4]} In this
chapter we shall catalog the different types of supersymmetry algebra
possible in higher dimensions, and use them to classify supermultiplets of
particles.
